% =============================================================================
\section{Discussion}
\label{sec:discussion}
% =============================================================================

\subsection{What a reported metric value does and does not mean}

The results in Section~\ref{sec:res_cohort} put a bound on how much a single
reported number can carry. Two studies can both report ``MMD$^2 = 0.15$ against
$500$ real BraTS slices'' while differing by a factor of two in the quantity
they intend to describe, purely through reference composition. Nothing in the
number, the sample size, or the estimator's unbiasedness reveals the
difference. Unbiasedness is a property of the estimator with respect to a fixed
pair of distributions; it says nothing about whether the empirical reference
represents the population it stands for.

This is the generative-evaluation counterpart of a problem that discriminative
medical imaging settled some time ago. Subject-level rather than image-level
splits are standard there, because slices from one subject leak between train
and test \citep{varoquaux2022machine}, and multi-site datasets are known to
carry recoverable scanner signatures \citep{glocker2019machine}. Generative
evaluation inherited neither convention, because it inherited its protocols from
natural-image benchmarks where neither problem exists. The reference set in a
generative study plays the structural role that a test split plays in a
discriminative one, and it deserves the same care.

\subsection{Why the conventional significance test is the wrong instrument}

Attaching a permutation test to a kernel two-sample statistic is technically
correct and practically uninformative. Its null --- that real and generated
samples come from the same distribution --- is false by construction for every
generative model anyone would submit for evaluation. Our measurements bear this
out: the test assigns $z > 340$ to an in-domain DDPM and to colour fundus
photography alike (Table~\ref{tab:cohortnull}), separating neither from the
other nor from anything else.

The behaviour under improving references is worse than merely uninformative.
Because the permutation null is built by shuffling labels within the pooled
sample, a more heterogeneous real set inflates the null variance and
\emph{depresses} the reported statistic. We say heterogeneous rather than
representative deliberately: we have no sampling frame for the target
population, so a wider subject draw is demonstrably more diverse but not
thereby shown to be more representative. A
practitioner who improves their reference from $17$ to $100$ subjects would see
the conventional significance of their result fall, from $z = 423$ to $z = 347$,
while the evaluation actually became better resolved. The inter-cohort
statistic moves the other way, from $15.7$ to $89.0$.

We do not claim the inter-cohort reference distribution is the only correct construction. It requires
subject labels, which not every released dataset provides; it is defined at a
particular cohort size, so two studies must match cohort size to be comparable;
and with $K$ blocks it rests on $K(K-1)/2$ non-independent pairs, so its tail
should be read as descriptive rather than as a calibrated $p$-value. We report
$z_{\mathrm{cohort}}$ and the ratio to the \emph{mean} observed cohort distance
for that reason, rather than leaning on $p_{\mathrm{cohort}}$, whose resolution
is bounded by the number of pairs, or on a ratio to the maximum, which is an
order statistic that is not comparable across conditions contributing different
numbers of pairs. What we do claim is that a null built from real
cohorts answers the question a practitioner is actually asking, and that the
permutation null does not.

\subsection{Reading depth deliberately}

The depth audit yields a practical result and a cautionary one. The practical
result is that a single frozen encoder supports several evaluation tasks at
different depths, at no additional cost: all twelve representations come from
one forward pass, so reading \Lb{8} for fidelity and \Lb{10} for near-duplicate
detection is free. The cautionary result is that the final block, which is what
the literature reads by default, is the best choice for none of the four tasks
we measured, and is $22\%$ below the fidelity plateau. We name \Lb{8} as the
maximum but not as ``the'' optimal depth: \Lb{7}, \Lb{8} and \Lb{9} are not
distinguishable from one another, and the defensible statement is that the
optimum lies in a mid-to-late plateau, not at the final block.

We are explicit that our depth recommendations are empirical rather than
principled. We do not have a theory that predicts why fidelity discriminability
should peak at \Lb{8} in this encoder, and we would not expect the specific
number to transfer to a different backbone or a different modality. What we
expect to transfer is the methodology: the depth profile is cheap to measure,
the figures of merit above do not presume an answer, and any study using a
foundation model as a metric backbone can run the audit for its own setting in
minutes.

The memorization reversal described in Section~\ref{sec:res_depth} is worth
restating as a general warning. On a \textsc{head} reference we previously
concluded that shallow blocks were better than mid-depth blocks for
near-duplicate detection. On a subject-disjoint reference the ordering reverses
and becomes monotone in the opposite direction. The reference protocol did not
shift a number by a few percent; it inverted a methodological conclusion. Any
depth study, ablation, or backbone comparison conducted on a convenience
reference is exposed to the same risk.

\subsection{Limitations}

\paragraph{One dataset for the cohort analysis}
The cohort-heterogeneity results are measured on BraTS, whose multi-site
composition and subject-ordered filenames make the effect easy to isolate. We
expect the phenomenon wherever a reference is assembled from a multi-site
collection of volumetric studies, but we have not measured its magnitude
elsewhere, and the specific inflation factors reported here should not be
assumed to transfer. Whether datasets curated from a single site show a
correspondingly smaller effect is an open question that the protocol in
Section~\ref{sec:cohortnull} is designed to answer.

\paragraph{One encoder for the depth analysis}
The depth grid uses RadioDINO-s16. The finding that depth is task-dependent is
likely general; the specific optima are not claimed to be.

\paragraph{Generators are not state of the art}
The DDPM and WDM-3D samples serve as objects of measurement, not as
contributions. A stronger generator would sit closer to the cohort null and
would make the resolution argument of Section~\ref{sec:res_null} more pressing,
not less: the closer a generator comes to real data, the more the reference
protocol determines whether the remaining gap can be detected at all.

\paragraph{Lesion-level specificity is a partial result}
Using expert glioma masks with area-matched mirror and texture-matched
controls, the RadioDINO MMD responds $1.4$--$1.8\times$ more strongly to
perturbation confined to tumour tissue than to matched healthy tissue, where an
Inception-based MMD responds at $1.05$--$1.08\times$. The effect is measurable
only in the upper part of the noise sweep: bootstrap intervals exclude unity at
$\sigma = 60$ ($1.57$, $[1.03, 1.97]$) and $\sigma = 90$ ($1.39$,
$[1.06, 1.71]$) but not at $\sigma = 45$ ($1.71$, $[0.94, 2.14]$), and at
$\sigma \le 30$ the ratio is numerically degenerate because the control-region
response is near zero. We report it as suggestive support for the domain
backbone rather than as an established property.

\paragraph{Coverage}
Our coverage result is a negative one about a specific estimator in a specific
embedding. It does not establish that no coverage diagnostic can work here, and
we did not evaluate alternatives such as $\alpha$-precision/$\beta$-recall
\citep{alaa2022faithful} at every depth. The recommendation is to run the
mode-drop validation before reporting any coverage number, not to abandon
coverage as a concept.

% =============================================================================
\section{Recommended protocol}
\label{sec:protocol}
% =============================================================================

The following is what we would ask of a paper reporting a distributional metric
for medical image generation, and what our released code implements by default.

\begin{enumerate}[leftmargin=*, topsep=2pt, itemsep=3pt]
\item \textbf{Report the number of distinct subjects behind every image set,
      not only the number of images.} A reference of $500$ slices from $17$
      subjects and one from $400$ are different instruments. This costs one
      number in a table and makes the remaining recommendations checkable.
\item \textbf{Build references by subject, not by filename.} Sample subjects
      first, then draw slices balanced across them and shuffled within them.
      Sorted-order convenience selects a site.
\item \textbf{Keep real-versus-real comparisons subject-disjoint.} A null built
      by splitting slices within subjects measures near-duplicate distance and
      is optimistically small.
\item \textbf{Contextualise against inter-cohort variation.} Report
      $z_{\mathrm{cohort}}$ and the ratio to the \emph{mean} real-versus-real
      cohort distance, at a stated cohort size, alongside the raw score. Avoid
      ratios to the maximum: it is an order statistic whose expectation grows
      with the number of cohort pairs compared. Report the permutation
      $p$-value if desired, but do not read it as evidence of generator
      quality.
\item \textbf{Audit read-out depth; do not copy ours.} Run the depth audit
      for your own encoder --- it costs one forward pass. Our optimum
      (\Lb{7}--\Lb{9} for fidelity, with those three not distinguishable from
      one another) is a property of this encoder and this data, and should not
      be transplanted.
\item \textbf{Validate every diagnostic before reporting it.} Mode drop must
      reduce measured coverage; injected duplicates must be recovered at the
      injected rate. A diagnostic that fails its own validation should not
      appear in a results table.
\item \textbf{Use one implementation per baseline metric.} FID computed from
      two different Inception variants differs by tens of points on identical
      images.
\end{enumerate}

% =============================================================================
\section{Conclusion}
\label{sec:conclusion}
% =============================================================================

We set out to audit a domain-pretrained embedding for generative medical image
evaluation and found that the more consequential variable lies upstream of the
embedding entirely. On BraTS, how the real reference cohort is assembled changes
the reported discrepancy by up to a factor of two, moves the real-versus-real
null by more than it moves the real-versus-generated score, and --- in the
worst case that the most convenient protocol happens to select --- produces a
distance between two sets of real images that exceeds the distance between real
and generated ones. Reference composition is not a second-order detail in this
setting; it is comparable in magnitude to the effect being measured.

We proposed calibrating against a null built from distances between disjoint
real patient cohorts, which asks whether a generator differs from the reference
by more than routine inter-cohort variation. Unlike the conventional
permutation test, which assigns overwhelming significance to every generator and
becomes \emph{less} significant as the reference improves, the cohort-calibrated
statistic tracks reference quality in the right direction and yields an
interpretable effect size.

Within that protocol, a depth-resolved audit of RadioDINO-s16 shows that
encoder depth is task-dependent and that the final block --- the default
read-out --- is optimal for none of the four tasks measured, falling $22\%$
below the fidelity plateau at \Lb{7}--\Lb{9}, whose members are not
distinguishable from one another. One widely used diagnostic, $k$-NN
precision/recall coverage, fails its own mode-drop validation at every depth
through neighbourhood-radius inflation, and we recommend against reporting it
in this feature space without validation. For per-image scoring, the choice of
distance matters more than depth.

None of these findings requires a new metric. They ask for a more careful
account of what an existing one is measured against.

% =============================================================================
\section*{Code and data availability}
% =============================================================================

Code implementing the metric, the inter-cohort reference distribution, the
depth audit, the literature audit record and
every experiment reported here is available at
\url{https://github.com/Nishan-Charlie/m3-cohort-audit}. BraTS 2021
\citep{baid2021rsna} and the LGG collection \citep{buda2019association} are
available from their respective repositories under their own terms; we
redistribute no patient data. Scripts to regenerate every table and figure from
raw data are included, together with the exact subject lists defining each
cohort block.

\section*{CRediT authorship contribution statement}

\textbf{Sathiyamohan Nishankar:} Conceptualization, Methodology, Software,
Validation, Formal analysis, Investigation, Data curation, Writing --- original
draft, Writing --- review \& editing, Visualization.

\section*{Declaration of competing interest}

The author declares no known competing financial interests or personal
relationships that could have appeared to influence the work reported in this
paper.
